% -*- mode: latex -*- %

\providecommand{\lipconst}{\mathsf{L}}
\newcommand{\uinit}{u_n^{\textup{init}}}
\newcommand{\uerror}{\epsilon}

\section{Semi-parametric approaches}
\label{sec:semi-param}

The preceding analysis highlights distinctions between a
likelihood-based approach using all the labels, as
in~\eqref{eq:maximum-likelihood}, and the robust but somewhat slower rates
that majority vote estimators enjoy, as in
Proposition~\ref{prop:lowd-misspcfd-majority-vote-m}.  That full-label
estimators' performance so strongly depends on the fidelity of the link
(recall Corollary~\ref{corollary:lowd-misspcfd-logistic-regression})
suggests that we target estimators achieving the best of both worlds: learn
both a link function (or collection thereof) and refit the model using
\emph{all} the labels. While we mainly focus on efficiency bounds for
procedures we view as closer to practice---fitting a model based on
available data---in this section, we provide a few results targeting more
efficient estimation schemes through semiparametric estimation approaches.

As developing full semi-parametric efficiency bounds would
lengthen and dilute the paper, we develop a few relatively general and
simple convergence results into which we can essentially plug in
semiparametric estimators. In distinction from standard results in
semiparametric theory (e.g.~\cite[Ch.~25]{VanDerVaart98}
or~\cite{BickelKlRiWe98}), the results require little more than consistent
estimation of the links $\sigma_j\opt$ to recover $1/m$ (optimal) scaling in
the asymptotic covariance, as the special structure of our classification
problem allows more nuanced calculations; we assume each labeler (link
function) generates labels for each of the $n$ datapoints $X_1, \ldots,
X_n$, but we could relax the assumption at the expense of
cumbersome notation.  We give two example applications of the general
theory: the first (Sec.~\ref{sec:semiparametric}) analyzing a full pipeline
for a single index model setting, which robustly estimates direction $u\opt$,
the link $\sigma\opt$, and then re-estimates $\theta\opt$; the second
assuming a stylized black-box crowdsourcing mechanism that provides
estimates of labeler reliability, highlighting how even in some
crowdsourcing scenarios, there could be substantial advantages to using full
label information.

\subsection{Master results}
\label{sec:master-results-sem}

For our specializations, we first provide master results that
allow semi-parametric estimation of the link functions.  We consider
Lipschitz symmetric link functions. Define for $\lipconst > 0$
\begin{align*}
  \fsymlinkL := \{\sigma \mid \sigma \not \equiv 1/2\text{ is non-decreasing,
    symmetric, and} ~ \lipconst\text{-Lipschitz continuous} \}
  \subset \fsymlink.
\end{align*}
We consider the general case where there are $m$ distinct labeler link
functions $\sigma_1\opt, \ldots, \sigma_m\opt$. To eliminate ambiguity
in the links, we
normalize the model with $\ltwo{\theta\opt} = 1$ so $\theta\opt
= u\opt$. We write $\vec{\sigma} = (\sigma_1, \ldots, \sigma_m)$ to distinguish from the typical case. For
$(x, y) \in \R^d \times \{\pm 1\}^m$ define
\begin{equation*}
  \loss_{\vec{\sigma}, \theta}(y \mid x)
  \defeq \frac{1}{m} \sum_{j = 1}^m \loss_{\sigma_j,\theta}(y_j \mid x),
\end{equation*}
which allows us to consider both the standard margin-based loss and the case
in which we learn per-labeler measures of quality.
We can then define the population
loss
\begin{equation*}
	L(\theta, \vec{\sigma})
	\defeq \frac{1}{m} \sum_{j = 1}^m
	\E[\loss_{\sigma_j, \theta}(Y_j \mid X)].
\end{equation*}
For any sequence $\{\vec{\sigma}_n\} \subset (\fsymlinkL)^m$ of (estimated)
links
and data
$(X_i, Y_i)$ for $Y_i = (Y_{i1}, \ldots, Y_{im})$, we define the semi-parametric
estimator
\begin{align*}
  \spe \defeq \argmin_{\theta \in \R^d}
  \frac{1}{n} \sum_{i = 1}^n
  \loss_{\vec{\sigma}_n, \theta}(Y_i \mid X_i).
\end{align*}

We will demonstrate both consistency and
asymptotic normality under appropriate convergence and regularity
assumptions for the link functions. We assume there is a (semiparametric)
collection $\fsymlinksp \subset (\fsymlinkL)^m$ of link functions of
interest, which may coincide with $(\fsymlinkL)^m$ but may be smaller,
making estimation easier.  Define the distance $d_{\fsymlinksp}$ on $\R^d
\times \fsymlinksp$ by
\begin{align}
  d_{\fsymlinksp}\left((\theta_1, \vec{\sigma}_1), (\theta_2, \vec{\sigma}_2)
  \right)
  \defeq
  \ltwo{\theta_1 - \theta_2}
  + \LtwoPbold{\vec{\sigma}_1(-Y\<X, u\opt\>)
    - \vec{\sigma}_2(-Y\<X, u\opt\>)}.
  \label{eqn:link-distance}
\end{align}
We make the following assumption.
\begin{assumption}
  \label{assumption:sp-general}
  The links $\vec{\sigma}\opt \in \fsymlinksp$ are normalized so
  that $\P(Y_j = y \mid X = x) = \sigma\opt_j(y \<x, u\opt\>)$, and
  the sequence $\{\vec{\sigma}_n\} \subset \fsymlinksp$ is consistent:
  \begin{equation*}
    \LtwoPbold{\vec{\sigma}_n(-Y\<X, u\opt\>)
      - \vec{\sigma}\opt(-Y\<X, u\opt\>)}
    \cp 0.
  \end{equation*}
  Additionally, the mapping $(\theta, \vec{\sigma}) \mapsto
  \nabla_\theta^2 \poploss(\theta,
  \vec{\sigma})$ is continuous for $d_{\fsymlinksp}$ at $(u\opt,
  \vec{\sigma}\opt)$.
\end{assumption}
The continuity of $\nabla_\theta^2 \poploss(\theta, \vec{\sigma})$ at
$(u\opt, \vec{\sigma}\opt)$ allows us to develop local asymptotic normality.
To see that we may expect the assumption to hold, we give reasonably simple
conditions sufficient for it, including that the collection of links
$\fsymlinksp$ is sufficiently smooth or the data distribution is continuous
enough. (See Appendix~\ref{proof:d-continuity-guarantees} for a proof.)

\begin{lemma} \label{lem:d-continuity-guarantees}
  Let $d_{\fsymlinksp}$ be the distance~\eqref{eqn:link-distance}. Let
  Assumption~\ref{assumption:x-decomposition} hold, where $|Z| > 0$ with
  probability one, has nonzero and continuously differentiable density
  $p(z)$ on $(0, \infty)$ satisfying $\lim_{z \to s} z^2 p(z) = 0$ for $s
  \in \{0, \infty\}$. The mapping $(\theta, \vec{\sigma}) \mapsto
  \nabla_\theta^2 \poploss(\theta, \vec{\sigma})$ is continuous for
  $d_{\fsymlinksp}$ at $(u\opt, \vec{\sigma}\opt)$ whenever $\Ep [\ltwo{X}^4] <
  \infty$ and either of the following conditions holds:
  \begin{enumerate}[leftmargin=2em,label=(C.\arabic*)]
  \item For any  $\vec{\sigma} = (\sigma_1, \ldots, \sigma_m) \in \fsymlinksp$, $\sigma_j'$ are Lipschitz continuous.
  \item $X$ has continuous density on $\R^d$.
  \end{enumerate}
\end{lemma}
We can now present the master result for semi-parametric
approaches, which characterizes the asymptotic behavior of the
semi-parametric estimator with the variance function
\begin{align*}
  C_{m, \vec{\sigma}\opt} := \frac{1}{m} \cdot
  \frac{\frac{1}{m} \sum_{j = 1}^m \E[\sigma_j\opt(Z) (1 - \sigma_j\opt(Z))]}{
    (\frac{1}{m} \sum_{j = 1}^m
    \E[{\sigma_j\opt}'(Z)])^2}.
\end{align*} 
\begin{theorem}
  \label{theorem:semiparametric-master}
  Let Assumption~\ref{assumption:x-decomposition} hold and assume $|Z| >
  0$ with probability one and has nonzero and continuous density $p(z)$ on
  $(0, \infty)$.  Let Assumption~\ref{assumption:sp-general} hold and assume
  that $\Ep [\ltwo{X}^4] < \infty$.  Then $\sqrt{n}(\spe - u\opt)$ is
  asymptotically normal, and the normalized estimator $\what{u}\sem_{n,m} =
  \spe / \ltwos{\spe}$ satisfies
	\begin{equation*}
		\sqrt{n} \prn{\what{u}\sem_{n,m} - u\opt}
		\cd \normal\prn{0, 	C_{m, \vec{\sigma}\opt}
			\prn{\proj_{u\opt}^\perp \Sigma \proj_{u\opt}^\perp}^\dagger}.
	\end{equation*}
\end{theorem}
\noindent
See Appendix~\ref{proof:semiparametric-master} for proof.
Notably, Theorem~\ref{theorem:semiparametric-master} exhibits optimal
$1/m$ covariance scaling whenever $\E[{\sigma\opt_j}'(Z)] \gtrsim 1$.


\subsection{A single index model}
\label{sec:semiparametric}

Our first example application of Theorem~\ref{theorem:semiparametric-master}
is to a single index model. We present a multi-phase estimator that first
estimates the direction $u\opt = \theta\opt / \ltwos{\theta\opt}$, then uses
this estimate to find a (consistent) estimate of the link $\sigma\opt$,
which we can then substitute directly into
Theorem~\ref{theorem:semiparametric-master}.
We defer all proofs of this section to Appendix~\ref{sec:proof-nonparametric},
which also includes a few auxiliary results that we use to prove the
results proper.

We present the the abstract convergence result for link functions first,
considering a scenario where we have an initial guess $\uinit$ of the
direction $u\opt$, independent of $(X_i, Y_{ij})_{i \le n,j\le m}$, for
example constructed via a small held-out subset of the data. We set
\begin{equation*}
  \vec{\sigma}_n = \argmin_{\vec{\sigma} \in \fsymlinksp}
  \sum_{i = 1}^n\sum_{j = 1}^m
  \left(\sigma_j(\<\uinit, X_i\>) - Y_{ij}\right)^2,
\end{equation*}
where $\fsymlinksp \subset (\fsymlinkL)^m$ and so it consists of
nondecreasing $\lipconst$-Lipschitz link functions with $\sigma(0) = \half$.
We assume that for all $n$, there exists a (potentially random) $\epsilon_n$
such that
\begin{equation*}
  \ltwos{\uinit - u\opt} \le \uerror_n.
\end{equation*}
\begin{proposition}
  \label{proposition:consistency-of-sigma}
  Let $X_i$ be vectors with $\E[\ltwo{X}^k] < \infty$, where
  $k \ge 2$. Then with probability $1$,
  there is a finite (random)
  $C < \infty$ such that for all large enough $n$,
  \begin{equation*}
    \LtwoPbold{\vec{\sigma}_n(Y\<u\opt, X\>)
      - \vec{\sigma}\opt(Y\<u\opt, X\>)}^2
    % = \LtwoQ{\sigma_n - \sigma\subopt}^2
    \le C \left[
      n^{\frac{2}{3k} - \frac{2}{3}}
      + \uerror_n^2
      + \lipconst n^{-\frac{k}{2(k + 1)}}
      \right].
  \end{equation*}
\end{proposition}
\noindent
The proof is more or less a consequence of standard convergence results
for nonparametric function estimation, though we include
it for completeness in Appendix~\ref{sec:proof-consistency-of-sigma}
as it includes a few additional technicalities because
of the initial estimate of $u\opt$.

Summarizing, we see that a natural procedure is available: if we have models
powerful enough to accurately estimate the conditional label probabilities
$Y \mid X$, then Proposition~\ref{proposition:consistency-of-sigma} coupled
with Theorem~\ref{theorem:semiparametric-master} shows that we can achieve
estimation with near-optimal asymptotic covariance. In particular, if
$\uinit$ is consistent (so $\uerror_n \cp 0$), then $\spe$ induces a
normalized estimator $\what{u}\sem_n = \spe / \ltwos{\spe}$ satisfying
$\sqrt{n}(\what{u}\sem_n - u\opt) \cd \normal(0, C_{m,\vec{\sigma}}
\prns{\proj_{u\opt}^\perp \Sigma \proj_{u\opt}^\perp}^\dagger)$.



\subsection{Crowdsourcing model}
\label{sec:crowdsourcing}

%\ha{the results in this section need more motivation. Currently we only compare to majority-vote estimator however the main point of the works we reference in this section is that using this model we should do better than majority-vote aggregation? Can we provide some comparison in rates to what they do? Or to any aggregation strategy?}.

Crowdsourcing typically targets estimating rater reliability, then using
these reliability estimates to recover ground truth labels as accurately as
possible, with versions of this approach central since at least
\citeauthor{DawidSk79}'s Expectation-Maximization-based
approaches~\cite{DawidSk79, WhitehillWuBeMoRu09, RaykarYuZhVaFlBoMo10}.  We
focus here on a simple model of rater reliability,
highlighting how---at least in our stylized model of classifier
learning---by combining a crowdsourcing reliability model and still using
all labels in estimating a classifier, we can achieve asymptotically
efficient estimates of $\theta\opt$, rather than the robust but slower
estimates $\what{\theta}_{n,m}\mv$ arising from ``cleaned'' labels.

%% To relate our results, we consider the case where the annotators are
%% heterogeneous. We still consider a binary classification model with multiple
%% labelers, that is our data consists of $\left\{(X_i, Y_{i1}), \dots, (X_i,
%% Y_{im})\right\}_{i=1}^n$, and $m$ is the number of annotators. Instead of
%% assuming $Y_{ij} \mid X_i$ are identically distributed for $j=1,\dots, m$,
%% we only assume that they are independent, allowing distinct link functions.

We adopt \citeauthor{WhitehillWuBeMoRu09}'s roughly ``low-rank'' model for
label generation~\citep{WhitehillWuBeMoRu09}: for binary classification with
$m$ labelers and distinct link functions $\sigma_i\opt$, model the
difficulty of $X_i$ by $\beta_i \in (-\infty, \infty)$, where
$\sgn(\beta_i)$ denotes the true class $X_i$ belongs to.  A parameter
$\alpha_j$ models the expertise of annotator $j$, and the probability
labeler $j$ correctly classifies $X_i$ is
\begin{align*}
  \Prb(Y_{ij} = 1) = \frac{1}{1 +\exp(-\alpha_i \beta_j)}.
\end{align*}
(See also \citet{RaykarYuZhVaFlBoMo10}.) 
The focus in these papers was to construct gold-standard labels
and datasets $(X_i, Y_i)$; here, we take the alternative perspective
we have so far advocated to show how using all labels can yield strong
performance.

%\ha{This is only true for majority vote with ERM. It would be interesting
%to show this for any aggregation strategy. Unfortunately this isn't the
%case currently.}

We thus adopt a semiparametric approach: we model the labelers, assuming a
black-box crowdsourcing model that can infer each labeler's ability, then
fit the classifier.  We represent labeler $j$'s expertise by a scalar
$\alpha_j\opt \in (0, \infty)$. Given data $X_i = X$ and the normalized
$\theta\opt = u\opt$, we assume a modified logistic link
\begin{align*}
  \Prb(Y_{ij} = 1 \mid X_i = x) = \frac{1}{1 + \exp (-\alpha\opt_j  \< \theta\opt, x \>)} = \sigma\lr (\alpha_j\opt  \< u \opt, x \>),
\end{align*}
so $\alpha_j\opt = \infty$ represents an omniscient labeler while
$\alpha_j\opt = 0$ means the labeler chooses random labels regardless of the
data. Let $\alpha\opt = (\alpha_1\opt, \ldots, \alpha_m\opt) \in
\R^m_{+}$. Then, in keeping with the plug-in approach of the master
Theorem~\ref{theorem:semiparametric-master}, we assume the blackbox
crowdsourcing model generates an estimate $\alpha_n = (\alpha_{n,1},
\ldots, \alpha_{n,m}) \in \R_+^m$ of $\alpha^\star$ from the data $\{(X_i,
(Y_{i1}, \ldots, Y_{im}))\}_{i=1}^n$.


%In other crowdsourcing literatures, \cc{Ref} developed other types of carefully designed algorithms including blablabla. The main theme that crowdsourcing community focuses on is assigning an as clean label as possible to data $i$ given the whole dataset by inferring the classification ability for each individual labeler.

%The paper however is purely empirical and lacks theoretical guarantees. Besides, another potential drawback for this modeling is that it assumes the annotator would make mistake at a constant probability $1-\alpha_j$ or $1-\beta_j$ even if the task is easy. Nonetheless, this paper shows that it's possible to learn a classifier and thus cover the learning procedure.

% There are more literatures to be discussed. Under this general setup, we discuss what is being calibrated and develop theory for training given a crowdsourcing black box.

%\subsubsection{What is a calibrated estimation?} 
%
%The first question might be a bit philosophical, but what is being calibrated when the annotators would have different predicting probabilities due to their different backgrounds (e.g. knowledge, expertise, or simply being lazy to label)? A short answer is, this still amounts to the same ``calibration error'' that we studied in the i.i.d. setup, that is $\ltwo{\what{\theta} - \theta\opt}$. To be specific, we would expect calibration in the following sense.
%
%\paragraph{Posterior calibration.} Suppose the expertise of each labeler $j$ is summarized in $\alpha_j$. It is reasonable to assume our labelers represent a certain group of people that we're interested in, that is $\alpha_j \simiid \mathbb{P}_\alpha$ we would want our machine learning procedure to produce a consistent estimate of the posterior probability, or the posterior difficulty of the data $X=x$. In this sense, being calibrated means
%\begin{align*}
%	\int \what{\Prb}\prn{Y = 1 \mid X =x, \alpha} d \Prb_\alpha \approx \int \Prb \prn{Y = 1 \mid X =x, \alpha} d \Prb_\alpha.
%\end{align*}
%
%\paragraph{Point-wise calibration.} Suppose for any new annotator with expertise $\alpha = t$, it holds uniformly in $t$ that
%\begin{align*}
%	\what{\Prb}(Y = 1 \mid X = x, \alpha = t) \approx \Prb (Y = 1 \mid X = x, \alpha = t),
%\end{align*}
%we say our learning procedure is point-wise calibrated. This allows generalization not only in the data domain $X \in \sX$ but also in the population domain $\alpha \in \mathcal{A}$. To be specific, imagine we have another model that is trained to infer expertise $\what{\alpha}$ for a new group of people, the plug in estimator $\what{\Prb}(Y = 1 \mid X = x, \what{\alpha})$ can be well-calibrated for this group of population. A more realistic application would be selected censoring for sensitive images that might cause different levels of disturbance for different people groups. 


%\subsubsection{Calibrated estimation using a blackbox crowdsourcing algorithm}
%Suppose $m$ and $d$ are fixed.
% It is clear that being posterior calibrated or point-wise calibrated both amount to $\ltwo{\what{\theta} - \theta\opt} \to 0$.

We consider the algorithm using the blackbox crowdsourcing model and
empirical risk minimization with the margin-based loss
$\loss_{\vec{\sigma}_n, \theta}$ as in Section~\ref{sec:master-results-sem},
with
$\vec{\sigma}_n(t)
\defeq (\sigma\lr(\alpha_{n,1}t),\dots, \sigma\lr(\alpha_{n,m}t))$,
or equivalently using the rescaled logistic loss
\begin{align*}
  \loss_{\vec{\sigma}_n, \theta}(y \mid x)
  = \frac{1}{m} \sum_{j=1}^m \loss\lr_\theta (y_j \mid \alpha_{n,j} x) .
\end{align*}
This allows us to apply our general semiparametric
Theorem~\ref{theorem:semiparametric-master} as long as the crowdsourcing
model produces a consistent estimate $\alpha_n \cp \alpha^\star$ (see
Appendix~\ref{proof:crowd-sourcing-estimator} for a proof):
\begin{proposition}
  \label{prop:crowd-sourcing-estimator}
  Let Assumption~\ref{assumption:x-decomposition} hold, $|Z| > 0$ have
  nonzero and continuous density $p(z)$ on $(0, \infty)$, and $\Ep
  [\ltwo{X}^4] < \infty$. If $\alpha_n \cp \alpha\opt \in \R^m$, then
  $\sqrt{n}(\spe - u\opt)$ is asymptotically normal, and the normalized
  estimator $\what{u}\sem_{n,m} = \spe / \ltwos{\spe}$ satisfies
  \begin{equation*}
    \sqrt{n} \prn{\what{u}\sem_{n,m} - u\opt}
    \cd \normal\prn{0,
      \frac{1}{\sum_{j = 1}^m \E[\sigma\lr(\alpha_j\opt Z)
          (1 - \sigma\lr(\alpha_j\opt Z))]}
      \prn{\proj_{u\opt}^\perp \Sigma \proj_{u\opt}^\perp}^\dagger}.
  \end{equation*}
%	\cc{If $|\alpha_j\opt| \geq \epsilon$ (i.e. are constant away from being just pure noise), then $C_m(t) \asymp 1/m$}
\end{proposition}

By Proposition~\ref{prop:crowd-sourcing-estimator}, the semiparametric
estimator $\what{u}\sem_{n,m}$ is efficient when the rater reliability
estimates $\alpha_n \in \R^m$ are consistent.  It is also immediate that if
$\alpha_j\opt \le \alpha_{\max} < \infty$ are bounded, then the asymptotic
covariance multiplier $C_{m,\alpha\opt} = (\sum_{j = 1}^m
\E[\sigma\lr(\alpha_j\opt Z) (1 - \sigma\lr(\alpha_j\opt Z))])^{-1} =
O(1/m)$, so we recover the $1/m$ scaling of the MLE, as opposed to the
slower rates of the majority vote estimators in
Section~\ref{sec:robustness-model-misspecification}.  At this point, the
refrain is perhaps unsurprising: using all the label information can yield
much stronger convergence guarantees.
